% Topic presentation -- 20 minutes. No animations or paper screenshots.
\documentclass[aspectratio=169,11pt]{beamer}
\usepackage[utf8]{inputenc}
\usepackage[T1]{fontenc}
\usepackage[english]{babel}
\usepackage{amsmath,amssymb}
\usepackage{booktabs}
\usetheme{default}
\usecolortheme{seahorse}
\setbeamertemplate{navigation symbols}{}
\setbeamertemplate{footline}[frame number]
\setbeamerfont{frametitle}{size=\large}
\setbeamercolor{alerted text}{fg=blue!65!black}

\newcommand{\source}[1]{\vfill{\tiny\color{black!60}#1}}

\title{AI-Assisted Learning and the Reversal\\of Long-Run Skill Degradation}
\subtitle{Topic presentation --- Track A}
\author{\texttt{wbgradost}}
\date{October 2026}

\begin{document}

\begin{frame}[plain]
  \titlepage
  \begin{center}\small
    \url{https://github.com/wbgradost/ai-project}
  \end{center}
\end{frame}

% ---------------------------------------------------------------------------
\section{Question and track}

\begin{frame}{1 \textperiodcentered{} The question}
  \begin{center}
    \Large Can AI assistance raise long-run skill when it makes the
    \alert{remaining human effort more effective for learning}?
  \end{center}
  \vspace{1em}
  \begin{columns}[T]
    \column{0.45\textwidth}
      \centering
      \textbf{Current production}\\[0.3em]
      AI substitutes for effort
    \column{0.1\textwidth}
      \centering\Large $\Longleftrightarrow$
    \column{0.45\textwidth}
      \centering
      \textbf{Skill acquisition}\\[0.3em]
      AI may complement effort
  \end{columns}
\end{frame}

\begin{frame}{1 \textperiodcentered{} Track A and target}
  \textbf{Baseline:} Aouad, Lykouris, and Zhong (2026)
  \begin{itemize}
    \item Proposition 3.5, p.~8: more AI shifts stationary skill downward.
  \end{itemize}
  \vspace{0.6em}
  \textbf{One changed assumption:}
  \[
    \lambda(e) \quad\longrightarrow\quad
    \lambda(e,a)=\lambda_0+\kappa e(1+\eta a).
  \]
  \vspace{0.3em}
  \alert{Target:} an exact two-state condition for a \emph{local} reversal,
  plus the threshold at which effort crowd-out regains dominance.
  \source{Aouad, Lykouris, and Zhong, arXiv:2605.11350v1; conference version EC 2026.}
\end{frame}

% ---------------------------------------------------------------------------
\section{Baseline and changed assumption}

\begin{frame}{2 \textperiodcentered{} Baseline worker problem}
  A myopic worker chooses effort for current utility:
  \[
    \max_{e\geq0}\; p(s+e+a)-\gamma e,
    \qquad \gamma>0.
  \]
  \begin{itemize}
    \item $p$: increasing, concave; $s,e,a$ are perfect substitutes.
    \item Learning does not enter the worker's current objective.
  \end{itemize}
  \vspace{0.4em}
  Kuhn--Tucker conditions:
  \[
    e^*\geq0,\quad \gamma-p'(s+e^*+a)\geq0,\quad
    e^*[\gamma-p'(s+e^*+a)]=0.
  \]
\end{frame}

\begin{frame}{2 \textperiodcentered{} Optimal effort: one-for-one crowd-out}
  Let
  \[
    x^*=\max\arg\max_{x\geq0}\{p(x)-\gamma x\}.
  \]
  Then Proposition 2.1 gives
  \[
    \boxed{e^*(s,a)=(x^*-s-a)^+.}
  \]
  \begin{columns}[T]
    \column{0.48\textwidth}
      \textbf{If $s+a<x^*$}
      \[
        \frac{\partial e^*}{\partial a}=-1,
        \qquad s+e^*+a=x^*.
      \]
    \column{0.48\textwidth}
      \textbf{If $s+a\geq x^*$}
      \[
        e^*=0.
      \]
  \end{columns}
  \source{Baseline result: Aouad, Lykouris, and Zhong (2026), Proposition 2.1.}
\end{frame}

\begin{frame}{2 \textperiodcentered{} Baseline skill dynamics}
  Continuous-time birth--death chain:
  \[
    s_i \xrightarrow{\ \lambda(e^*(s_i,a))\ } s_{i+1},
    \qquad
    s_i \xrightarrow{\ \mu\ } s_{i-1}.
  \]
  Local balance implies
  \[
    \frac{\pi_{i+1}(a)}{\pi_i(a)}
      =\frac{\lambda(e^*(s_i,a))}{\mu}.
  \]
  \vspace{0.4em}
  Higher $a$ $\Rightarrow$ lower $e^*$ $\Rightarrow$ lower upward rates
  $\Rightarrow$ more stationary mass on low skill.
  \begin{block}{Proposition 3.5 (baseline)}
    For $a_\ell<a_h$, the distribution under $a_\ell$ first-order
    stochastically dominates the distribution under $a_h$.
  \end{block}
\end{frame}

\begin{frame}{2 \textperiodcentered{} The assumption I change}
  \[
    \boxed{\lambda(e,a)=\lambda_0+\kappa e(1+\eta a)}
  \]
  \[
    \lambda_0>0,\qquad \kappa>0,\qquad \eta\geq0.
  \]
  \begin{align*}
    \lambda_e &= \kappa(1+\eta a)>0, &
    \lambda_{ea} &= \kappa\eta\geq0.
  \end{align*}
  \begin{itemize}
    \item For fixed $a$: linear and weakly concave in effort.
    \item $\eta=0$: affine special case of the original technology.
    \item Worker problem unchanged: effort remains $(x^*-s-a)^+$.
  \end{itemize}
\end{frame}

% ---------------------------------------------------------------------------
\section{Expected result}

\begin{frame}{3 \textperiodcentered{} Two states: stationary distribution}
  Let $s_L<s_H$ and define
  \[
    \lambda_L(a)=\lambda(e^*(s_L,a),a).
  \]
  \begin{center}
    $s_L\ \underset{\mu}{\overset{\lambda_L(a)}{\rightleftarrows}}\ s_H$
  \end{center}
  Stationary flow balance:
  \[
    \pi_L\lambda_L=\pi_H\mu,
    \qquad \pi_L+\pi_H=1.
  \]
  Therefore
  \[
    \boxed{\pi_H(a)=\frac{\lambda_L(a)}{\lambda_L(a)+\mu}.}
  \]
  \alert{The sign of $d\pi_H/da$ is exactly the sign of
  $d\lambda_L/da$.}
\end{frame}

\begin{frame}{3 \textperiodcentered{} Competing effects of AI}
  Set $b=x^*-s_L>0$. While $0\leq a<b$,
  \[
    e_L^*(a)=b-a
  \]
  and
  \[
    \lambda_L(a)
      =\lambda_0+\kappa(b-a)(1+\eta a).
  \]
  Total effect:
  \begin{align*}
    \frac{d\lambda_L}{da}
      &= \underbrace{\kappa\eta e_L^*}_{\text{learning complementarity}}
       - \underbrace{\kappa(1+\eta a)}_{\text{effort crowd-out}}\\[0.3em]
      &= \kappa[\eta(b-2a)-1].
  \end{align*}
  \[
    \boxed{\frac{d\pi_H}{da}>0
      \quad\Longleftrightarrow\quad \eta(b-2a)>1.}
  \]
\end{frame}

\begin{frame}{3 \textperiodcentered{} Candidate proposition: two thresholds}
  \begin{columns}[T]
    \column{0.47\textwidth}
      \textbf{Weak complementarity}
      \[
        \eta b\leq1
      \]
      $\pi_H(a)$ is weakly decreasing: baseline degradation survives.
    \column{0.49\textwidth}
      \textbf{Strong complementarity}
      \[
        \eta b>1
      \]
      Initial reversal, ending at
      \[
        \boxed{a^\dagger=\frac{\eta b-1}{2\eta}.}
      \]
  \end{columns}
  \vspace{0.8em}
  \begin{center}
    $0$ \quad $\xrightarrow{\ \pi_H\ \uparrow\ }$ \quad
    $a^\dagger$ \quad $\xrightarrow{\ \pi_H\ \downarrow\ }$ \quad
    $b$ \quad $\xrightarrow{\ \pi_H\ \text{constant}\ }$
  \end{center}
  \small At $a=b$, effort is zero; learning returns to $\lambda_0$.
\end{frame}

\begin{frame}{3 \textperiodcentered{} Exact pairwise reversal condition}
  For two active-effort levels $0\leq a_\ell<a_h<b$,
  \begin{align*}
    \lambda_L(a_h)-\lambda_L(a_\ell)
      &=\kappa(a_h-a_\ell)
        [\eta(b-a_h-a_\ell)-1].
  \end{align*}
  Since $\pi_H=\lambda_L/(\lambda_L+\mu)$,
  \[
    \boxed{
    \pi_H(a_h)>\pi_H(a_\ell)
    \iff
    \eta(b-a_h-a_\ell)>1.}
  \]
  \begin{alertblock}{Scope}
    This reverses Proposition 3.5 only for qualifying pairs of assistance
    levels. If $a_h\geq b>a_\ell$, the higher level eliminates effort and the
    reversal cannot hold.
  \end{alertblock}
\end{frame}

% ---------------------------------------------------------------------------
\section{Why not already solved}

\begin{frame}{4 \textperiodcentered{} What the baseline appendices establish}
  \begin{itemize}
    \item Appendix A.1 derives the stationary birth--death distribution.
    \item Appendix A.4 proves Proposition 3.5 using
      $\lambda=\lambda(e)$ and monotonicity in effort.
    \item Appendix D allows convex effort costs, but the learning rate still
      has no direct AI argument.
    \item The current arXiv record contains v1; the authors list an EC 2026
      conference version.
  \end{itemize}
  \vspace{0.5em}
  \alert{Missing object:} no AI--effort cross-partial in the transition
  technology, hence no closed-form stationary reversal condition.
  \source{Aouad, Lykouris, and Zhong (2026), Appendices A and D.}
\end{frame}

\begin{frame}{4 \textperiodcentered{} Closest related result: Davies (2026)}
  \textbf{Course Design in the Age of AI}
  \begin{itemize}
    \item Myopic student chooses work or full delegation; teacher designs tasks.
    \item AI quality affects both the marginal product of effort and delegation.
    \item Theorem 4: higher quality raises learning for high-skill students but
      lowers it for low-skill students.
  \end{itemize}
  \vspace{0.4em}
  \begin{block}{What this project adds to the Aouad et al. model}
    Keep its worker problem and stationary Markov structure; place
    complementarity directly in $\lambda(e,a)$; derive a closed-form stationary
    reversal condition and assistance threshold.
  \end{block}
  \small No claim that AI--learning complementarity itself is new.
  \source{Davies, arXiv:2607.18735v3, Theorem 4; Shen--Tamkin (2026); Wu et al. (2026).}
\end{frame}

% ---------------------------------------------------------------------------
\section{Plan and risks}

\begin{frame}{5 \textperiodcentered{} Work plan after topic approval}
  \begin{enumerate}
    \item \textbf{Proof:} all two-state boundary/equality cases; then exact
      sufficient conditions for an $N$-state ordering.
    \item \textbf{Symbolic checks:} verify stationary identities, derivatives,
      and pairwise differences.
    \item \textbf{Numerics:} parameter grids around $a^\dagger$ and $b$;
      separate evidence from proof.
    \item \textbf{Lean:} formalize
      \begin{itemize}
        \item the two-state stationary share,
        \item the derivative-sign equivalence,
        \item the pairwise condition, including positivity and piecewise cases.
      \end{itemize}
  \end{enumerate}
  \vspace{0.3em}
  \centering
  \alert{Proof statement first; simulations and Lean only after it is fixed.}
\end{frame}

\begin{frame}{5 \textperiodcentered{} Main risks and disciplined fallbacks}
  \begin{itemize}
    \item \textbf{Different $N$-state thresholds.}
      State sufficient rate-by-rate conditions; retain the exact two-state result.
    \item \textbf{A local result is read as global.}
      Report both thresholds $a^\dagger$ and $b$, with every domain restriction.
    \item \textbf{Myopia omits the value of future skill.}
      Keep it for baseline comparability; state forward-looking choice as a limitation.
    \item \textbf{$\lambda(e,a)$ is not jointly concave.}
      Require concavity in effort conditional on $a$, which is all the result uses.
  \end{itemize}
  \vspace{0.4em}
  \centering
  \Large\alert{Question for discussion: is local reversal the right scope?}
\end{frame}

\end{document}
